% GLOB-GEO0007 v1.0.0. Geometry and projections follow the CIVITA library,
% notched_cylinder.png, and the geometry figure design standard.
% Build from this directory with SOURCE_DATE_EPOCH=0 and
% pdflatex -halt-on-error -no-shell-escape notched_cylinder.tex. Convert the
% PDF with pdftocairo -svg notched_cylinder.pdf notched_cylinder.svg.
%
% Definition: right circular cylinder x^2 + y^2 <= (diameter/2)^2,
% 0 <= z <= length, with O at the center of the reference end face z = 0 and
% the z axis on the cylinder axis. A circumferential notch is centered at
% z_c = length/2 and spans z_c -/+ notch_width/2 axially. Its depth is radial,
% from the outer surface to the deepest point, so the root radius is
% R_root = diameter/2 - notch_depth.
%
% Four independent parameters (length, diameter, notch_depth, notch_width).
% In an axial section the notch has parallel sides ending in a semicircle;
% notch_tip_radius = 0.5 * notch_width is a derived parameter of the library.
% notch_depth includes the semicircular tip, so the straight-sided radial
% portion is notch_depth - notch_width/2, which the library requires to be
% strictly positive.
%
% The reference draws +z downward; the design standard projects +z upward, so
% the same right-handed frame appears rotated, with z = 0 at the bottom.
\documentclass[border=1pt]{standalone}
\usepackage[T1]{fontenc}
\usepackage{lmodern}
\usepackage{tikz}
\usetikzlibrary{3d,calc}
% Shared formatting for four-view CIVITA geometry figures.
\usetikzlibrary{arrows.meta}

\definecolor{CivitaGeometry}{HTML}{242424}
\definecolor{CivitaHidden}{HTML}{555555}
\definecolor{CivitaCenterline}{HTML}{525252}
\definecolor{CivitaConstruction}{HTML}{737373}
\definecolor{CivitaAxes}{HTML}{245B83}
\definecolor{CivitaDimension}{HTML}{202020}
\definecolor{CivitaSurface}{HTML}{ECEFF1}
\definecolor{CivitaText}{HTML}{1B1B1B}
\colorlet{CivitaSideSurface}{CivitaSurface!92!CivitaGeometry}
\colorlet{CivitaEndSurface}{CivitaSurface!48!white}

\newcommand{\CivitaPanelWidth}{180}
\newcommand{\CivitaPanelHeight}{134}
\newcommand{\CivitaColumnGap}{12}
\newcommand{\CivitaRowGap}{14}
\newcommand{\CivitaPanelPadding}{8}

\tikzset{
  civita/figure/.style={x=1pt,y=1pt,line cap=round,line join=round},
  civita/edge/.style={draw=CivitaGeometry,line width=.7pt},
  civita/hidden/.style={draw=CivitaHidden,line width=.45pt,dash pattern=on 2pt off 2pt},
  civita/centerline/.style={draw=CivitaCenterline,line width=.45pt,
    dash pattern=on 4pt off 1.5pt on .6pt off 1.5pt},
  civita/construction/.style={draw=CivitaConstruction,line width=.35pt,dotted},
  CivitaArrow/.tip={Latex[length=4pt,width=3pt]},
  civita/dimension/.style={draw=CivitaDimension,line width=.45pt,
    arrows=CivitaArrow-CivitaArrow},
  civita/extension/.style={draw=CivitaDimension,line width=.35pt},
  civita/dimension extension/.style={civita/extension,shorten <=2pt,shorten >=-2pt},
  civita/axis/.style={draw=CivitaAxes,line width=.65pt,arrows=-CivitaArrow},
  civita/label/.style={text=CivitaText,font=\sffamily\fontsize{8}{9}\selectfont,inner sep=1pt},
  civita/axis label/.style={civita/label,text=CivitaAxes},
  civita/title/.style={text=CivitaText,font=\sffamily\bfseries\fontsize{9}{10}\selectfont,
    inner sep=0pt,anchor=north west},
  civita/front/.style={fill=CivitaSurface},
  civita/side/.style={fill=CivitaSideSurface},
  civita/end/.style={fill=CivitaEndSurface}
}

\newcommand{\CivitaCanvas}{%
  \pgfmathsetmacro{\CivitaFigureWidth}{2*\CivitaPanelWidth+\CivitaColumnGap}%
  \pgfmathsetmacro{\CivitaFigureHeight}{2*\CivitaPanelHeight+\CivitaRowGap}%
  \fill[white] (0,0) rectangle (\CivitaFigureWidth,\CivitaFigureHeight);%
}

% Panel indices follow the design standard, with row one at the top.
\newenvironment{CivitaPanel}[3]{%
  \pgfmathsetmacro{\CivitaPanelX}{(#2-1)*(\CivitaPanelWidth+\CivitaColumnGap)}%
  \pgfmathsetmacro{\CivitaPanelY}{(2-#1)*(\CivitaPanelHeight+\CivitaRowGap)}%
  \begin{scope}[shift={(\CivitaPanelX,\CivitaPanelY)}]%
    \node[civita/title] at (\CivitaPanelPadding,{\CivitaPanelHeight-\CivitaPanelPadding}) {#3};%
}{\end{scope}}

\newcommand{\CivitaOrigin}[3][]{%
  \fill[CivitaAxes] (#2) circle[radius=1.3pt];%
  \node[civita/axis label,anchor=#3,#1] at (#2) {$O$};%
}

\newcommand{\CivitaAxis}[4]{%
  \draw[civita/axis] (#1) -- (#2) node[civita/axis label,anchor=#3] {$#4$};%
}

% Named feature points, dimension level in panel coordinates, and symbolic label.
\newcommand{\CivitaHorizontalDimension}[4]{%
  \coordinate (civita-dimension-level) at (0,#3);%
  \draw[civita/dimension extension] (#1) -- (#1 |- civita-dimension-level);%
  \draw[civita/dimension extension] (#2) -- (#2 |- civita-dimension-level);%
  \draw[civita/dimension] (#1 |- civita-dimension-level) -- (#2 |- civita-dimension-level)
    node[civita/label,midway,below=1.5pt] {#4};%
}

% Named feature points, dimension column, label placement options, and label.
\newcommand{\CivitaVerticalDimension}[5]{%
  \coordinate (civita-dimension-column) at (#3,0);%
  \draw[civita/dimension extension] (#1) -- (civita-dimension-column |- #1);%
  \draw[civita/dimension extension] (#2) -- (civita-dimension-column |- #2);%
  \draw[civita/dimension] (civita-dimension-column |- #1) -- (civita-dimension-column |- #2)
    node[civita/label,midway,#4] {#5};%
}

\newcommand{\CivitaParameter}[1]{\textsf{\detokenize{#1}}}

\pdfinfoomitdate=1

% Illustrative construction dimensions in mm. All printed labels are symbolic.
\newcommand{\NCLength}{160}      % length
\newcommand{\NCDiameter}{80}     % diameter
\newcommand{\NCNotchDepth}{18}   % notch_depth (radial)
\newcommand{\NCNotchWidth}{10}   % notch_width (axial)

% Derived parameter of the library: notch_tip_radius = 0.5 * notch_width.
\pgfmathsetmacro{\NCTipRadius}{0.5*\NCNotchWidth}
% Derived construction values (not parameters of the library).
\pgfmathsetmacro{\NCR}{\NCDiameter/2}                   % outer radius
\pgfmathsetmacro{\NCRroot}{\NCR-\NCNotchDepth}          % root radius
\pgfmathsetmacro{\NCRt}{\NCRroot+\NCTipRadius}          % tip center radius
\pgfmathsetmacro{\NCzc}{\NCLength/2}                    % notch center
\pgfmathsetmacro{\NCzm}{\NCzc-\NCNotchWidth/2}          % lower notch wall
\pgfmathsetmacro{\NCzp}{\NCzc+\NCNotchWidth/2}          % upper notch wall

% Constraints of geometry_library.json: strict positivity of the four
% independent parameters and the three geometric constraints.
\newcommand{\NCRequire}[2]{%
  \pgfmathparse{(#1)?1:0}%
  \ifnum\pgfmathresult=0
    \PackageError{civita-notched-cylinder}{Violated constraint: #2}{Check the parameter values.}%
  \fi}
\NCRequire{\NCLength>0 && \NCDiameter>0 && \NCNotchDepth>0 && \NCNotchWidth>0}{length, diameter, notch_depth, notch_width > 0}
\NCRequire{\NCNotchDepth<0.5*\NCDiameter}{notch_depth < 0.5 * diameter (positive root radius)}
\NCRequire{\NCNotchWidth<\NCLength}{notch_width < length}
\NCRequire{\NCNotchDepth>0.5*\NCNotchWidth}{notch_depth > 0.5 * notch_width (positive straight-sided portion)}

% Cubic arc about the tip center, in the meridian half-plane with radial unit
% vector (#1,#2,0): center radius #3, polar angles #4 -> #5 measured in the
% (radial, z) plane, control length (4/3) tan((b-a)/4) r.
\newcommand{\NCArc}[5]{%
  ({#1*(#3+\NCTipRadius*cos(#4))},{#2*(#3+\NCTipRadius*cos(#4))},{\NCzc+\NCTipRadius*sin(#4)})
  .. controls
  ({#1*(#3+\NCTipRadius*(cos(#4)-4/3*tan((#5-#4)/4)*sin(#4)))},
   {#2*(#3+\NCTipRadius*(cos(#4)-4/3*tan((#5-#4)/4)*sin(#4)))},
   {\NCzc+\NCTipRadius*(sin(#4)+4/3*tan((#5-#4)/4)*cos(#4))})
  and
  ({#1*(#3+\NCTipRadius*(cos(#5)+4/3*tan((#5-#4)/4)*sin(#5)))},
   {#2*(#3+\NCTipRadius*(cos(#5)+4/3*tan((#5-#4)/4)*sin(#5)))},
   {\NCzc+\NCTipRadius*(sin(#5)-4/3*tan((#5-#4)/4)*cos(#5))})
  .. ({#1*(#3+\NCTipRadius*cos(#5))},{#2*(#3+\NCTipRadius*cos(#5))},{\NCzc+\NCTipRadius*sin(#5)})}
% Semicircular tip contour at radial direction (#1,#2), from the lower wall
% through the deepest point to the upper wall.
\newcommand{\NCTip}[2]{\NCArc{#1}{#2}{\NCRt}{270}{180} -- \NCArc{#1}{#2}{\NCRt}{180}{90}}
\newcommand{\NCTipReversed}[2]{\NCArc{#1}{#2}{\NCRt}{90}{180} -- \NCArc{#1}{#2}{\NCRt}{180}{270}}
% Axial section of the notched cylinder in the plane spanned by (#1,#2,0)
% and z: the meridian profile on both sides of the axis.
\newcommand{\NCOutline}[2]{%
  ({-#1*\NCR},{-#2*\NCR},0) -- ({#1*\NCR},{#2*\NCR},0) -- ({#1*\NCR},{#2*\NCR},\NCzm)
  -- \NCTip{#1}{#2} -- ({#1*\NCR},{#2*\NCR},\NCzp) -- ({#1*\NCR},{#2*\NCR},\NCLength)
  -- ({-#1*\NCR},{-#2*\NCR},\NCLength) -- ({-#1*\NCR},{-#2*\NCR},\NCzp)
  -- \NCTipReversed{-#1}{-#2} -- ({-#1*\NCR},{-#2*\NCR},\NCzm) -- cycle}
% Radius of the notch neck at height z (lower wall to upper wall).
\newcommand{\NCNeckRadius}[1]{\NCRt-sqrt(max(0,\NCTipRadius^2-(#1-\NCzc)^2))}

\colorlet{NCNotchSurface}{CivitaSurface!78!CivitaGeometry}

\newcommand{\NCOrthoScale}{.60}
\newcommand{\NCIsoScale}{.54}
\pgfmathsetmacro{\NCIsoHorizontal}{\NCIsoScale/sqrt(2)}
\pgfmathsetmacro{\NCIsoDown}{\NCIsoScale/sqrt(6)}
\pgfmathsetmacro{\NCIsoUp}{2*\NCIsoDown}
\newcommand{\NCAxisLength}{28}
\newcommand{\NCCenterOverrun}{10}
\pgfmathsetmacro{\NCIsoAxisLength}{1.4*\NCR}
\pgfmathsetmacro{\NCIsoCenterTop}{\NCLength+\NCR}
\newcommand{\NCNeckSamples}{24}

% Local points shared by the orthographic projections.
\newcommand{\NCPoints}{%
  \coordinate (nc-o) at (0,0,0);
  \coordinate (center-low) at (0,0,-\NCCenterOverrun);
  \coordinate (center-high) at (0,0,\NCLength+\NCCenterOverrun);
  \coordinate (axis-x) at (\NCAxisLength,0,0);
  \coordinate (axis-y) at (0,\NCAxisLength,0);
  \coordinate (axis-z) at (0,0,\NCAxisLength);
}

\begin{document}
\begin{tikzpicture}[civita/figure]
  \CivitaCanvas

  \begin{CivitaPanel}{1}{1}{(a) 3D View}
    % View along (-1,-1,-1), with positive z projected upward. A ray toward
    % the viewer always moves toward larger z, so material can only be
    % hidden by material above it. The specimen is therefore painted in
    % ascending z: lower cylinder, notch neck, upper cylinder. Circles are
    % drawn in their own x-y planes and projected by the same affine map;
    % the side contours of the neck lie in the meridian planes at polar
    % angles 135 and -45 degrees, where the projected circles reach their
    % horizontal extremes.
    \begin{scope}[shift={(90,27)},x={(-\NCIsoHorizontal pt,-\NCIsoDown pt)},
      y={(\NCIsoHorizontal pt,-\NCIsoDown pt)},z={(0pt,\NCIsoUp pt)}]
      \NCPoints
      % Lower cylinder, 0 <= z <= z_-.
      \path[civita/front] (0,0,0) circle[radius=\NCR];
      \path[civita/front] ({\NCR*cos(135)},{\NCR*sin(135)},0) -- ({\NCR*cos(-45)},{\NCR*sin(-45)},0)
        -- ({\NCR*cos(-45)},{\NCR*sin(-45)},\NCzm) -- ({\NCR*cos(135)},{\NCR*sin(135)},\NCzm) -- cycle;
      \begin{scope}[canvas is xy plane at z=0]
        \draw[civita/hidden] (135:\NCR) arc[start angle=135,end angle=315,radius=\NCR];
        \draw[civita/edge] (-45:\NCR) arc[start angle=-45,end angle=135,radius=\NCR];
      \end{scope}
      \draw[civita/edge] ({\NCR*cos(135)},{\NCR*sin(135)},0) -- ({\NCR*cos(135)},{\NCR*sin(135)},\NCzm)
        ({\NCR*cos(-45)},{\NCR*sin(-45)},0) -- ({\NCR*cos(-45)},{\NCR*sin(-45)},\NCzm);
      % Lower notch wall z = z_- (an annulus facing +z).
      \path[civita/end,civita/edge] (0,0,\NCzm) circle[radius=\NCR];
      % Notch neck, z_- <= z <= z_+: union of its projected circles, then its
      % two side contours (the semicircular tip seen edge-on).
      \foreach \i in {0,...,\NCNeckSamples}{
        \pgfmathsetmacro{\NCz}{\NCzm+\i*\NCNotchWidth/\NCNeckSamples}
        \pgfmathsetmacro{\NCr}{\NCNeckRadius{\NCz}}
        \path[fill=NCNotchSurface] (0,0,\NCz) circle[radius=\NCr];
      }
      \draw[civita/edge] \NCTip{-0.70711}{0.70711};
      \draw[civita/edge] \NCTip{0.70711}{-0.70711};
      % Upper cylinder, z_+ <= z <= length.
      \path[civita/front] (0,0,\NCzp) circle[radius=\NCR];
      \path[civita/front] ({\NCR*cos(135)},{\NCR*sin(135)},\NCzp) -- ({\NCR*cos(-45)},{\NCR*sin(-45)},\NCzp)
        -- ({\NCR*cos(-45)},{\NCR*sin(-45)},\NCLength) -- ({\NCR*cos(135)},{\NCR*sin(135)},\NCLength) -- cycle;
      \begin{scope}[canvas is xy plane at z=\NCzp]
        \draw[civita/edge] (-45:\NCR) arc[start angle=-45,end angle=135,radius=\NCR];
      \end{scope}
      \path[civita/end,civita/edge] (0,0,\NCLength) circle[radius=\NCR];
      \draw[civita/edge] ({\NCR*cos(135)},{\NCR*sin(135)},\NCzp) -- ({\NCR*cos(135)},{\NCR*sin(135)},\NCLength)
        ({\NCR*cos(-45)},{\NCR*sin(-45)},\NCzp) -- ({\NCR*cos(-45)},{\NCR*sin(-45)},\NCLength);
      % Hidden rear halves of the two notch rims, where the upper cylinder
      % covers them, so the groove reads as continuous around the axis.
      \begin{scope}
        % (band traced counterclockwise, like the circle, so the nonzero
        % rule yields their union)
        \clip (0,0,\NCzp) circle[radius=\NCR]
          ({\NCR*cos(135)},{\NCR*sin(135)},\NCzp) -- ({\NCR*cos(135)},{\NCR*sin(135)},\NCLength)
          -- ({\NCR*cos(-45)},{\NCR*sin(-45)},\NCLength) -- ({\NCR*cos(-45)},{\NCR*sin(-45)},\NCzp) -- cycle;
        \begin{scope}[canvas is xy plane at z=\NCzm]
          \draw[civita/hidden] (135:\NCR) arc[start angle=135,end angle=315,radius=\NCR];
        \end{scope}
        \begin{scope}[canvas is xy plane at z=\NCzp]
          \draw[civita/hidden] (135:\NCR) arc[start angle=135,end angle=315,radius=\NCR];
        \end{scope}
      \end{scope}
      \coordinate (iso-x) at (\NCIsoAxisLength,0,0);
      \coordinate (iso-y) at (0,\NCIsoAxisLength,0);
      \coordinate (iso-center-high) at (0,0,\NCIsoCenterTop);
      \draw[civita/centerline] (axis-z) -- (iso-center-high);
    \end{scope}
    \CivitaAxis{nc-o}{iso-x}{north east}{x}
    \CivitaAxis{nc-o}{iso-y}{north west}{y}
    \CivitaAxis{nc-o}{axis-z}{south west}{z}
    \CivitaOrigin[yshift=-3pt]{nc-o}{north}
  \end{CivitaPanel}

  \begin{CivitaPanel}{1}{2}{(b) x-y View}
    % View along -z onto the end face z = length. The notch root circle at
    % mid-length is hidden.
    \begin{scope}[shift={(90,70)},x={(\NCOrthoScale pt,0pt)},
      y={(0pt,\NCOrthoScale pt)},z={(0pt,0pt)}]
      \NCPoints
      \coordinate (xy-x-minus) at (-\NCR,0,0);
      \coordinate (xy-x-plus) at (\NCR,0,0);
      \draw[civita/centerline] ({-\NCR-\NCCenterOverrun},0) -- (-\NCR,0) (\NCR,0) -- ({\NCR+\NCCenterOverrun},0);
      \draw[civita/centerline] (0,{-\NCR-\NCCenterOverrun}) -- (0,-\NCR) (axis-y) -- (0,{\NCR+\NCCenterOverrun});
      \draw[civita/hidden] (nc-o) circle[radius=\NCRroot];
      \draw[civita/edge] (nc-o) circle[radius=\NCR];
    \end{scope}
    \CivitaHorizontalDimension{xy-x-minus}{xy-x-plus}{36}{\CivitaParameter{diameter}}
    \CivitaAxis{nc-o}{axis-x}{south west}{x}
    \CivitaAxis{nc-o}{axis-y}{west}{y}
    \CivitaOrigin{nc-o}{north east}
  \end{CivitaPanel}

  \begin{CivitaPanel}{2}{1}{(c) x-z View}
    % View along +y. Axial section profile on both sides of the axis; the
    % notch rims at z_- and z_+ project as full-width lines.
    \begin{scope}[shift={(90,16)},x={(\NCOrthoScale pt,0pt)},
      y={(0pt,0pt)},z={(0pt,\NCOrthoScale pt)}]
      \NCPoints
      \draw[civita/centerline] (center-low) -- (center-high);
      \draw[civita/edge] \NCOutline{1}{0};
      \draw[civita/edge] (-\NCR,0,\NCzm) -- (\NCR,0,\NCzm) (-\NCR,0,\NCzp) -- (\NCR,0,\NCzp);
      \coordinate (xz-root) at (\NCRroot,0,\NCzc);
      \coordinate (xz-surface) at (\NCR,0,\NCzc);
      \coordinate (xz-rim-lo) at (\NCR,0,\NCzm);
      \coordinate (xz-rim-hi) at (\NCR,0,\NCzp);
      \coordinate (xz-left-lo) at (-\NCR,0,\NCzm);
      \coordinate (xz-left-hi) at (-\NCR,0,\NCzp);
    \end{scope}
    % notch_width: axial, at the notch on the -x side; arrows outside.
    \coordinate (xz-w-col) at ([xshift=-8pt]xz-left-lo);
    \draw[civita/dimension extension] (xz-left-lo) -- (xz-w-col |- xz-left-lo);
    \draw[civita/dimension extension] (xz-left-hi) -- (xz-w-col |- xz-left-hi);
    \draw[civita/extension] ([yshift=-7pt]xz-w-col |- xz-left-lo) -- ([yshift=7pt]xz-w-col |- xz-left-hi);
    \draw[civita/dimension,arrows=-CivitaArrow] ([yshift=-7pt]xz-w-col |- xz-left-lo) -- (xz-w-col |- xz-left-lo);
    \draw[civita/dimension,arrows=-CivitaArrow] ([yshift=7pt]xz-w-col |- xz-left-hi) -- (xz-w-col |- xz-left-hi);
    \node[civita/label,anchor=east] at ([xshift=-2pt]$(xz-w-col |- xz-left-lo)!.5!(xz-w-col |- xz-left-hi)$)
      {\CivitaParameter{notch_width}};
    % notch_depth: radial, from the original outer surface (thin reference
    % line across the opening) to the deepest point of the tip.
    \draw[civita/extension] ([yshift=-2pt]xz-rim-lo) -- ([yshift=2pt]xz-rim-hi);
    \draw[civita/dimension] (xz-root) -- (xz-surface);
    \node[civita/label,anchor=west] at ([xshift=2pt]xz-surface) {\CivitaParameter{notch_depth}};
    \CivitaAxis{nc-o}{axis-x}{south west}{x}
    \CivitaAxis{nc-o}{axis-z}{west}{z}
    \CivitaOrigin[xshift=-1.5pt]{nc-o}{north east}
  \end{CivitaPanel}

  \begin{CivitaPanel}{2}{2}{(d) y-z View}
    % View along -x. Same axisymmetric profile, same scale as (c).
    \begin{scope}[shift={(84,16)},x={(0pt,0pt)},
      y={(\NCOrthoScale pt,0pt)},z={(0pt,\NCOrthoScale pt)}]
      \NCPoints
      \draw[civita/centerline] (center-low) -- (center-high);
      \draw[civita/edge] \NCOutline{0}{1};
      \draw[civita/edge] (0,-\NCR,\NCzm) -- (0,\NCR,\NCzm) (0,-\NCR,\NCzp) -- (0,\NCR,\NCzp);
      \coordinate (yz-left) at (0,-\NCR,0);
      \coordinate (yz-left-top) at (0,-\NCR,\NCLength);
      \coordinate (yz-tip) at (0,\NCRroot,\NCzc);
      \coordinate (yz-out) at (0,\NCR,\NCzc);
    \end{scope}
    \CivitaVerticalDimension{yz-left}{yz-left-top}{47}{left=1pt,rotate=90,anchor=south}{\CivitaParameter{length}}
    % Derived tip radius: single leader through the tip center to the arc,
    % annotated with the library relation in smaller, secondary type.
    \draw[civita/dimension,arrows=-CivitaArrow] ([xshift=6pt]yz-out) -- (yz-tip);
    \node[civita/label,font=\sffamily\fontsize{7}{8}\selectfont,align=left,anchor=west]
      at ([xshift=6.5pt]yz-out)
      {\CivitaParameter{notch_width} =\\ 2\,\CivitaParameter{notch_tip_radius}};
    \CivitaAxis{nc-o}{axis-y}{south west}{y}
    \CivitaAxis{nc-o}{axis-z}{west}{z}
    \CivitaOrigin[xshift=-1.5pt]{nc-o}{north east}
  \end{CivitaPanel}
\end{tikzpicture}
\end{document}
